\scanpage{28}
% loc 1–3, continuation at upper edge
structure de $\unsure{\widetilde A|\check A}$-module sur $M$ tel que pour \unsure{tt|tout} $x\in M$, $\xi\longmapsto\xi x$ soit continue (i.e. s'annule sur un des $J_{\unsure{i|1|\prime}}\subset\check A\ \ldotp\ldotp\ldotp\ldotp\ \unsure{]|\exists}$. 
% loc 3–6, De =, on voit
De $=$, on voit que la donnée \deleted{\unsure{s}}d'une structure de coalgèbre \inserted{(\unsure{ass. unit})} sur le $k$-module libre $A$ revient à la donnée d'une \deleted{opération} structure de $k$-algèbre (ass. unitaire) sur $\check A$ qui soit \uline{continue}
% loc 7–8, car si; struck line and formula
[car si $\overset{\inserted{\unsure{k^{I}}}}{\overset{\scriptscriptstyle\unsure{\tikz[baseline=-.4ex]{\node[rotate=90,inner sep=0pt]{\(\scriptscriptstyle\simeq\)};}|\Vert}}{P}},\overset{\inserted{\unsure{k^{J}}}}{\overset{\tikz[baseline=-.4ex]{\node[rotate=90,inner sep=0pt]{\(\scriptscriptstyle\simeq\)};}}{Q}},\overset{\inserted{\unsure{k^{\Lambda}}}}{\overset{\scriptscriptstyle\Vert}{R}}$ sont libres, \deleted{les \unsure{accouplements|applications} \inserted{bil.} continus $\check P\times\check Q\unsure{\longrightarrow|-}\check R=\underset{\leftarrow}{\lim}_{\unsure{\alpha|\lambda}}\deleted{\unsure{R_{I}}}R_{\unsure{\Lambda_{\alpha}|\Lambda_2}}$.}
% loc 9–16, multiline Bil and Hom display
\[
\begin{aligned}
\unsure{\mathrm{bil}|\mathrm{Bil}}\,\mathrm{cont}_{k}\,\check P,\check Q;\check R)
 &\simeq \mathrm{Bil}\,\mathrm{cont}_{k}(\check P,\check Q;\underset{\leftarrow}{\lim}_{\gamma}(\check R_{\gamma}))\\
 &\simeq \underset{\leftarrow}{\lim}_{\gamma}\,[\mathrm{bil}\,\mathrm{cont}_{k}(\check P,\check Q;\check R_{\gamma})]\\
 &\simeq \underset{\leftarrow}{\lim}_{\gamma}\,\bigl[\underset{\unsure{\rightarrow|-}}{\lim}_{\alpha,\beta}\,
 \underbrace{\mathrm{Bil}_{k}(\check P_{\alpha},\check Q_{\beta};\check R_{\gamma})}_{\mathrm{Hom}_{k}(\check P_{\alpha}\otimes\check Q_{\beta};\check R_{\gamma})\,]}\bigr]\\
 &\deleted{\unsure{\simeq}}\\
 &\simeq \underset{\leftarrow}{\lim}_{\gamma}\,\underset{\rightarrow}{\lim}_{\alpha,\beta}\,\mathrm{Hom}_{k}(R_{\gamma},P_{\alpha}\otimes Q_{\beta})\\
 &\simeq \underset{\leftarrow}{\lim}_{\gamma}\,\mathrm{Hom}_{k}(R_{\gamma},P\otimes Q)\\
 &\simeq \mathrm{Hom}_{k}(R,P\otimes Q)\qquad !\quad ]
\end{aligned}
\]

% loc 17–21, On est dans
\deleted{On a également le même résultat} \inserted{On est dans le cas envisagé} si $A$ est limite \deleted{\unsure{projective}}\unsure{inductive} \uline{dénombrable} de sous-modules projectifs \inserted{de t.f.}, à morphismes de transition \deleted{ayant} des \deleted{suppléments} mon. \unsure{univ.|inv.} [i.e. ayant des inverses à g.] Mais dans le cas général, on ne peut oublier $V(A)$ \inserted{et ses opérations} au profit de $\check A=V(A)(k)$ $\ldots$

% loc 22–27, Module quotient and quotient map
Module quotient et sous-module stable par une co-opération de la co-algèbre $A$ : Soit $M\overset{\unsure{g|f}}{\unsure{\longrightarrow|-}}N'$ un ép., alors $V(N)\subset V(M)$ mono, et il est clair ce que signifie que $V(N)$ est stable ; on dira alors que $N$ est stable par la coop. de $N$. \unsure{Cel|Cela} signifie aussi que $\exists$ un hom $N'\overset{\omega_{N'}}{\longrightarrow}N'\otimes A$ rendant comm.
% loc 27–28, commutative square
\diagram{\begin{tikzcd}[ampersand replacement=\&,column sep=large]
M\arrow[r,"\omega_M"]\arrow[d] \& M\otimes A\arrow[d,"g\otimes\mathrm{id}_{A}"]\\
N'\arrow[r,"\omega_{N'}"] \& N'\otimes A
\end{tikzcd}}
% loc 29–30, last two lines
et cet $\omega_{N}$ est néc. unique et satisfait aux axiomes d'unité et associativité. Si $A$ est \unsure{un gp.|une alg.|var.} affine
% archive sheet number: 27
